% --------------------------------------------------
% LLNL proposal Template
% --------------------------------------------------
\documentclass[11pt]{article}
\usepackage[letterpaper, margin=1in]{geometry}
\usepackage{times}

% --------------------------------------------------
% Colors
% --------------------------------------------------
\usepackage[table]{xcolor}
\usepackage{tikz}

\definecolor{secblue}{HTML}{3B5E7F}
\definecolor{subsecblue}{HTML}{4F81BD}
\definecolor{hdrgray}{HTML}{A7B7C7}

\definecolor{boxblue}{HTML}{587590}
\definecolor{boxgray}{HTML}{E0E6E6}

% --------------------------------------------------
% Draft watermark
% --------------------------------------------------
%\usepackage[colorspec=.9]{draftwatermark}
%\SetWatermarkText{DRAFT}
%\SetWatermarkScale{1}

% --------------------------------------------------
% Extra packages
% --------------------------------------------------
\usepackage[subfigure]{tocloft}
\usepackage{inconsolata}
\usepackage{abstract}
\usepackage{subfigure}
\usepackage{graphicx}
\usepackage{hyperref}

\hypersetup{
  colorlinks,
  citecolor=violet,
  linkcolor=blue,
  urlcolor=blue
}
\usepackage{wrapfig}
\usepackage{url}
\usepackage{pdfpages}
\usepackage{ctable}
\usepackage{multirow}
\usepackage[titletoc]{appendix}
\usepackage{lipsum}
\usepackage{mdframed}
\usepackage{fancyvrb}
\usepackage{tabu}
\usepackage{tabularx}
\usepackage{multirow}
\usepackage{xspace}

\usepackage{pdfpages}

% Sort multiple citations by citation number
\usepackage[sort]{natbib}

% Space reduction
\usepackage{caption}
\DeclareCaptionFont{subsecblue}{\color{subsecblue}}
\captionsetup{labelfont={small,subsecblue,bf},font=small}

\usepackage{enumitem}
\setlist{noitemsep,topsep=2pt}

% --------------------------------------------------
% Section heading format
% --------------------------------------------------
\usepackage{titlesec}

\titlespacing{\parameterization}{0em}{.4em}{.5em}
\titlespacing{\section}{0em}{.4em}{.1em}
\titlespacing{\subsection}{0em}{.4em}{.1em}
\titlespacing{\subsubsection}{0em}{.4em}{.1em}

\def\secformat{\sffamily\bfseries}

\titleformat{\section}{\large\secformat\color{secblue}}{\thesection}{1em}{}
\titleformat{\subsection}{\large\secformat\color{subsecblue}}{\thesubsection}{1em}{}
\titleformat{\subsubsection}{\normalsize\secformat\color{subsecblue}}{\thesubsubsection}{1em}{}
\titleformat{\paragraph}[runin]{\normalsize\secformat\color{subsecblue}}{\theparagraph}{1em}{}

\newcommand{\tablesec}[1]{%
        \vspace{.1em}
        \textbf{\color{subsecblue} #1}\newline}
% --------------------------------------------------
% Paragraphs
% --------------------------------------------------
\setlength{\parindent}{0em}
\setlength{\parskip}{.5em}


% --------------------------------------------------
% Make TOC fonts match sections
% --------------------------------------------------
\renewcommand\abstractnamefont{\normalsize\secformat\color{secblue}}

\renewcommand\contentsname{Table of Contents}

\renewcommand\cfttoctitlefont{\Large\secformat\color{secblue}}
\renewcommand\cftsecfont{\normalsize\secformat\color{secblue}}
\renewcommand\cftsubsecfont{\normalsize\sffamily\color{subsecblue}}
\renewcommand\cftsubsubsecfont{\normalsize\sffamily\color{subsecblue}}

\renewcommand\appendixpagename{\Large\secformat\color{secblue} Appendices}
\renewcommand\appendixtocname{\Large\secformat\color{secblue} Appendices}


% --------------------------------------------------
% Headers & Geometry
% --------------------------------------------------

% Header for other pages
\usepackage{fancyhdr}
\renewcommand{\headrulewidth}{0pt}
\pagestyle{fancy}
\lhead{
        \begin{tikzpicture}
                \node[white,draw=white,shape=rectangle,fill=hdrgray] at (-8, 12)
                {\sf\footnotesize \hbox to \textwidth{
                        Strong Foundations for Sustainability
                        \hfill
                        FY23 ASCR Software Sustainability Proposal}
                };
        \end{tikzpicture}
}
\chead{}
\rhead{}

% --------------------------------------------------
% Title, Author, Proposal info table
% --------------------------------------------------
\def\lbox{\cellcolor{boxblue}\color{white}\small\sffamily\bfseries}
\def\rbox{\cellcolor{boxgray}\small}
\taburulecolor{white}
%
% Title
\title{\vspace{-3em}\sf\huge\color{secblue}%
Strong Foundations for Sustainability}
%
% Author
%\author{\vspace{-1em}
%  Todd Gamblin, \href{mailto:tgamblin@llnl.gov}{tgamblin@llnl.gov}
%  Tzanio Kolev, \href{mailto:tzanio@llnl.gov}{tzanio@llnl.gov}
%  Ann Almgren, \href{mailto:asalmgren@lbl.gov}{asalmgren@lbl.gov}
%}
\date{\vspace{-4em}\today\vspace{-3em}}

% --------------------------------------------------
% Custom comamnds
% --------------------------------------------------
\newcommand{\ncite}{{\color{red}[cite]}}

% --------------------------------------------------
% Main document
% --------------------------------------------------
\begin{document}
\setcounter{page}{1}

\maketitle

\section{Vision}

%The Exascale Computing Project (ECP)
ASCR has created a broad ecosystem of software that
enables scientists across government, academia and industry to realize the full
potential of exascale machines in advancing the DOE mission.
%Like its biological namesake,
This ecosystem is thriving and interconnected,
%Researchers and developers constantly evolve the ECP
%software stack to adapt it for new needs and new mission applications. More than ever
%before (in DOE), the push to exascale has enabled {\it integration} across the stack,
%and {\it robust}, {\it shared} capabilities are leveraged by many consumer packages. The
%ECP stack provides many unique capabilities, but these capabilities cannot be sustained
%unless the software is maintained.
%but it is only
but it only continues to be
valuable if users have confidence in its reliability and longevity.
%This requires ongoing contributions, and a
A successful sustainment strategy requires growth: of the software itself,
but also of the stakeholder community and of the contributor base.
%should aim to grow the contributor base as large as possible.
We propose to achieve this through these key thrusts:

\begin{enumerate}

\item {\bf Core sustainment and integration.} We will create a project structure to sustain key
    software packages, including a critical mass of the most impactful software developed
    by ASCR-funded researchers.  We will ensure sustained integration of the components of the
    software stack through ongoing efforts in Spack.

\item {\bf Stakeholder growth.} We will aggressively lower barriers for adoption, attracting new
    users and contributors to maximize the value of DOE's investment and broaden our
    impact.  We will support multiple models for interactions with industry,
    enabling non-DOE funding to support software as appropriate.

\item {\bf Workforce development.} We will proactively seek to broaden and diversify the
    ASCR software community in multiple ways: through increasing the participation in
    existing outreach activities\footnote{See, e.g., the existing LF initiatives we can build on:
  \href{https://www.linuxfoundation.org/about/diversity-inclusivity}{https://www.linuxfoundation.org/about/diversity-inclusivity}};
    and through directly offering
    mentoring and internships with participating software teams.  We will partner with the
    Sustainable Horizons Institute, allowing us to maximize the impact of our efforts.
\end{enumerate}

%Spack, Kokkos, MFEM, and other ECP projects have grown
%  communities of hundreds of contributors, attracting the interest of industry just as
%  HPC has begun to take off in the cloud\footnote{Current studies project an 8\%
%  Compound Annual Growth Rate (CAGR) for cloud HPC.}. Now is the time to leverage this
%  confidence.

\section{Software Foundations}

One successful model of software sustainment outside of DOE is through software
foundations like the Apache Software Foundation (ASF), NumFOCUS, and the Linux
Foundation (LF), which have successfully scaled and sustained some of the largest open
source communities. Within LF, the Cloud Native Computing Foundation (CNCF) organizes a
consortium of independent projects around containers and Kubernetes. Importantly, a key
goal of CNCF is to create a {\it portable} software stack to prevent cloud vendor
lock-in, much as ECP has created a portable stack that spans GPUs. CNCF's projects are
widely supported across clouds and even garner contributions and support from cloud
vendors.

There is currently no analog of CNCF for HPC, and we aim to fill this gap. The ultimate
aim of this project is to create an HPC Foundation (HPCF) that can:
\begin{enumerate}
\item Serve as a neutral home for key projects in the HPC open source ecosystem, and
      provide a home for project infrastructure, e.g. Continuous Integration (CI).
\item Provide a transparent governance model that allows DOE, facilities, industry, and
      academia to sit at the same table and steer the direction of member projects, and
      allows member projects to accept donations and funds from foundation members to
      help sustain projects;
%\item Understand HPC trends, coordinating with CNCF and others to increase adoption and
%      integration;
\item Help develop the workforce of tomorrow by engaging them early in their careers.
\end{enumerate}

\section{Governance}

Neutral, transparent governance is critical for growing the contributor base for the ECP
stack, as it allows the broadest possible set of parties to be stakeholders. Currently,
all governance sits within DOE. While collaboration is encouraged, there are barriers
for industry collaborators who want to partake in roadmap planning. Shared ownership
eliminates barriers, allows competitors (like cloud vendors) to collaborate, and
instills confidence that the project will endure beyond any one organization's support.

We will thus adopt a model similar to
\href{https://github.com/kubernetes/steering/blob/main/charter.md}{that of Kubernetes},
where a top-level steering committee guides the project's strategic direction, bylaws,
and finances, but delegates all technical responsibilities to sub-groups. Initially, the
steering committee will comprise senior staff from ECP projects, facilities,
laboratories, ASCR, NNSA, DOE applications, and industry partners. This model is
familiar to industry, and it lowers barriers for external funding. Many companies are
already set up to sponsor similarly structured organizations.

% \section{Impact, Metrics, and Core Sustainment}

To present a familiar technical governance model, the steering committee will be
responsible for the ramp-up and ramp-down of Special Interest Groups (SIGs). SIGs are
groups of collaborators working together on particular key technical topics, which could
range from core feature development to a larger group of related projects focusing on
refactoring or integration efforts. SIGs can be funded or unfunded, but key projects
designated for core sustainment will be funded based on decisions by the steering
committee.

The steering committee will determine which projects to sustain by considering a number
of key metrics on a regular schedule. We propose the following guidelines now, but part
of our efforts in FY23 will be to work with the steering committee to establish guidelines
for FY24:

\begin{enumerate}
\item Broad usage (number of users, contributors, etc.) and potential for community
  growth;
\item Critical dependencies and capabilities for DOE exascale applications and software
  products;
\item Market trends and potential to establish industry, facility, or academic
  partnerships;
\item Potential to enhance the diversity of the larger HPC community.
  and
\item Key infrastructure for the sustainment and integration of the HPCF stack.
\end{enumerate}

The steering committee will evaluate metrics and projects on a yearly basis, and funded
projects will be required to establish yearly milestones.
%Milestone descriptions should
%be short (no more than a page or a single slide) but should demonstrate progress towards
%the objectives above. Milestones will be presented to the steering committee or by a
%designated subcommittee each year.
Inclusion will be determined by the above metrics, but
funding amounts for approved projects will be determined based on technical need.

\section{Growing SIGs into Foundation Projects}

We expect each SIG to have its own open meetings, to publicly document how external
parties can participate, and to include new potential partners and customers in the
project planning. We believe that enthusiastic users eventually become contributors, and
funded projects will be given assistance by the HPCF with project tooling such as
automation, hosting, establishing CI, interfacing with HPC facilities, and cultivating
industry partnerships. This will lower barriers for project developers and will
encourage further growth. We will establish a model similar to
\href{https://www.cncf.io/projects/}{CNCF's ``sandbox'', ``incubating'' and
  ``graduated'' projects} to grow project maturity and move projects to
production-readiness, and we will aggressively promote production-ready projects to the
broader HPC community. SIGs are not {\it necessarily} tied to funding; they can simply
be collaborations that arise in the course of open source development. We expect
projects to first show promise by our key metrics before being funded by DOE, but our
model opens the door to complementary funding.

The SIG structure provides a lightweight governance model for smaller projects, but it
allows larger projects to grow their own governance model in the hope that they could
become foundation projects in their own right. If a SIG begins to grow its own community
that needs more extensive structure and bylaws, it can become a separate foundation
projects within the HPCF. The HPCF and its steering committee will work to maintain
awareness of all member projects, establishing collaborations but not dictating their
governance models or development workflows. We do not believe there is a single best
model for technical decisions.

\section{Pilot Foundation Projects in FY23}

We cannot grow an overarching HPC foundation overnight, but we can prototype this
governance structure in FY23.
We are already in conversations with the
Linux Foundation. We will need to document project governance models and establish
procedures for establishing SIGs. We have interest from Amazon Web Services (AWS),
Google Cloud Platform (GCP), HPE, AMD, Kitware, and others in sponsoring an initial
foundation. In FY23 we will:

\begin{enumerate}
\item Establish an initial steering committee and SIG structure;
\item Establish an initial set of SIGs to cover a large fraction of the ECP software
  ecosystem; and
\item Establish Spack (and potentially others) as initial LF projects with dues and
  clear membership levels for industry, facility, labs, and other partners;
\item Begin discussions with LF about an overarching HPC foundation;
\item Work with ASCR to develop a model for DOE staff to work on funded SIGs within
  foundations;
\item Begin outreach to attract industry and academic partners to join us. We will explore potential partnership with the HPC4EI program in this area.
\end{enumerate}

Spack itself does not encompass the entire ECP software stack, but Spack has high
adoption and visibility in industry, and all ECP projects are deployed using Spack.
Given its centrality to the ECP software ecosystem, we aim to coordinate integration
among projects through Spack, just as E4S has done within ECP.
%Spack is the glue that
%holds the ecosystem together, and it allows us to present a unified external interface
%for deployment at HPC sites or in the cloud.

\section{Initial Special Interest Groups (SIGs)}

As mentioned above, SIGs are technical focus areas with publicly documented activities
and governance. In our view, SIGs can be funded if necessary, based on decisions by the
steering committee. We will work with steering committees to establish an initial set of
SIGs representing critical components of the ECP stack. We aim to cover as much of the
ECP stack as possible. Unlike SDKs, the SIGs do not need aim to produce a common
software product, but rather advance the field, similar to the co-design centers in the
ECP. Our aim is not to force projects to work together, but to find the engagement areas
that will inspire them to do so. The SIG model can accommodate different funding levels,
by varying the number of funded SIGs.

In FY23 we will stand up SIGs for Spack, Kokkos, MFEM/CEED, AMReX and other
projects with significant contributor bases.
These SIGs will be organized around topic areas such as: Software Stack Integration, Structured and
Unstructured Discretizations, Sparse Solvers and Time Integrators,
Programming Models and Compilers, Workflow, Visualization, Performance Tools, System Software and
IO, and Machine Learning (ML). This list is evolving rapidly, but it corresponds to SDKs
and co-design centers within ECP, as well as key collaboration directions identified by
ECP projects. Our goal is to coordinate with other funded sustainability projects to
cover projects that wish to engage with us on our foundation model. If certain projects
are funded through another sustainability project, they can still join and contribute to
a SIG here, but we should be able to fund the entire ECP stack if possible and
necessary.

\section{FY23 and Beyond}

In FY23, we will socialize our ideas with the community and incorporate changes based on
feedback. We will organize the steering committee and initial SIGs, then provide several
options to ASCR depending on the funding level. After FY23, we will work with LF to
establish the umbrella HPCF project.

Scientific computing is at an inflection point. Accelerated AI is ubiquitous, and
industry and cloud providers are beginning to look at accelerated computing for
HPC\footnote{The projected Compound Annual Growth Rate (CAGR) for cloud HPC is 8\%.}.
Today's exascale technologies are tomorrow's cloud offerings.
%A well maintained DOE stack can be the foundation of an accelerated cloud HPC ecosystem.
%and we can use our position of software leadership to shape the future of cloud and industry HPC.
%
With sufficient momentum, DOE software can shape the direction of industry HPC.
DOE open source has value that no company wishes to reinvent. A well established
software ecosystem ensures that we can take advantage of post-exascale hardware, and
blaze a path to a portable cloud HPC stack.




%\newpage
% --------------------------------------------------
% References
% --------------------------------------------------
%\nocite{*}
%\begingroup
%\section{References Cited}
%\renewcommand{\section}[2]{}
%\bibliographystyle{abbrv}
%{
%  \footnotesize
%  \setlength{\parskip}{0em}
%  \bibliography{ascr-sustainability-fy23}
%}
%\endgroup


% --------------------------------------------------
% Biographical Sketch
% --------------------------------------------------

% A biographical sketch (limited to two pages) is required for each PI,
% co-investigator, student, postdoc, and/or university collaborator. The
% information must be provided in the order and format specified in the
% LDRD Biographical Sketch Template: Do not submit personal information
% such as home address; home telephone, fax, or cell phone numbers; home
% e-mail address; date of birth; citizenship; driver’s license numbers;
% marital status; personal hobbies; etc.

%\newpage
%\includepdf[pages=-]{grant-cv.pdf}
%\section{Biographical Sketch}

%\includepdf[pages=-]{bio-sketches/gamblin.pdf}

\end{document}
